\documentclass[10pt,a4paper]{amsart}
\usepackage[margin=19mm]{geometry}
\usepackage{amsmath,amssymb,mathtools}
\usepackage[scr=rsfs,cal=euler]{mathalfa}
\usepackage{xfrac}
\usepackage[unicode,hidelinks,bookmarks=false]{hyperref}
\newcommand{\ie}{i.e.\ }
\newcommand{\cf}{cf.\ }
\newcommand{\resp}{respectively\ }
\newcommand{\Subsection}{\subsection}
\providecommand{\nobreakdash}{\nobreak\hbox{-}}
\setlength{\emergencystretch}{2em}
\multlinegap0pt
\usepackage{amssymb}
\usepackage{tikz}
\newtheorem{lemma}{Lemma}
\newtheorem{remark}{Remark}
\newtheorem{theorem}{Theorem}
\newcommand{\ZZ}{{\mathbb{Z}}}
\newcommand{\bc}{\begin{center}}
\newcommand{\ec}{\end{center}}
\renewcommand{\geq}{\geqslant}
\newcommand{\la}{\langle}
\newcommand{\ra}{\rangle}
\newcommand{\sss}{\sigma}
\title{On the trivial units property and the unique product property}
\author{Heiko Dietrich, Melissa Lee, Andre Nies, Marc Vinyals}
\date{Source: arXiv:2603.22640v2 (CC BY 4.0)}
\begin{document}
\maketitle

\noindent\emph{Source and licence.} On the trivial units property and the unique product property. Version arXiv:2603.22640v2 (CC BY 4.0), distributed by arXiv under the Creative Commons Attribution 4.0 International licence, http://creativecommons.org/licenses/by/4.0/, as recorded in the arXiv licence metadata for this exact version and shown on the abstract page https://arxiv.org/abs/2603.22640v2. Full licence notice: \texttt{licenses/arxiv-2603.22640-CC-BY-4.0.txt}.\par\smallskip

\noindent\textit{Editorial scope.} Counted unit: the article's theorem Hn\_solv\_wp (source0.tex lines 581--581). The article's complete original proof of it is delivered with it (source0.tex lines 583--596, one \texttt{proof} environment of 14 lines). No mathematics is written, repaired or re-derived: the statement and the proof are the article's own lines, in the article's own notation, and no retained line is substituted or re-flowed.  Retained uncounted as the source-local context the proof needs: lines 551--552; lines 555--557; lines 569--570; lines 574--575.  Nothing was omitted from any retained span, so the package carries no omission record.  No retained line contains a citation, so the wrapper carries no bibliography and no bibliography entry.  No retained line contains a figure environment, an image-inclusion command or drawing code, so no image is delivered and the package declares no asset dependency.  Editorial formatting only, in addition: an AMS \texttt{amsart} preamble that restates the article's own declarations and macros used by the retained lines, namely the environments \texttt{lemma}, \texttt{remark}, \texttt{theorem} on separate counters, together with the standard packages those lines need.  The retained environments print in this document as Lemma 1, Remark 1, Theorem 1 in order of appearance, whereas the article prints the same items with the numbering of its own full document.  The complete native source of the article as distributed in the arXiv e-print for this exact version is bundled unchanged as \texttt{sources/197032/source0.tex}.
\begin{align}\label{eq_fo} g\in K_n\text{ has finite order} \iff g^h\in \langle x_i \rangle \text{ for some $i$ and $h\in K_n$}.
\end{align}

\begin{lemma} \label{lem Kn}
There is an algorithm that can rewrite  every element in $K_n$ into a unique normal form  $x_{i_1}\ldots x_{i_k}x_n^z$ where successive $i_u$ and $i_{u+1}$ are distinct, $i_k\ne n$, and $z$ is some integer. 
\end{lemma}

\begin{remark} \label{Malcev} If the word problem for $G$ is solvable, the group is computable in the usual sense of computable algebra: there is a bijection  $\theta \colon G \to \mathbb N$ such that the images under $\theta$ of group products are computable. 
 \end{remark}

\begin{lemma} \label{lem: rec presented} Let $H$ be a recursively presented group and let   $w\in Z(H)$  be a non-trivial central element. If $L= H/\la w \ra$ has a solvable word problem, then so does~$H$.
\end{lemma}

  \begin{theorem}\label{Hn_solv_wp} The group $H_n$ has solvable word problem for $n\geq 4$. \end{theorem}

  \begin{proof} Define $L_n = H_n/\la w_n \ra$. In $L_n$, each $x_i^2=1$, and so the relator $w_n=1$ can be written as  $x_1 \ldots x_r =  x_{n}^{-1} \ldots   x_{r+1}^{-1}=  x_{n} \ldots   x_{r+1}$ for any $r\in\{2,\ldots,n-2\}$. Thus, we have a presentation
  \bc $L_n = \la\; x_1,  \dots , x_n \mid x_1^2,\;\ldots,\;x_n^2,\; x_1 \ldots x_r =  x_{n} \ldots   x_{r+1} \;\ra$,  \ec
which shows that $L_n\cong A*_\ZZ B$ is an amalgamated free product where 
 %%%%%%%
 \bc $A= \la\; x_1, \ldots, x_r \mid x_1^2,\ldots,x_r^2 \;\ra \quad\text{and}\quad B= \la\; x_{r+1}, \ldots, x_n \mid x_{r+1}^2,\ldots,x_n^2\; \ra$, \ec
 %%%%%%%
 and $\ZZ$ is embedded into $A $ via $1 \mapsto x_1 \ldots x_r$ and into $B$ via $1 \mapsto   x_{n} \ldots x_{r+1} $. Indeed, since $r$ lies in $\{2,\ldots, n-2\}$, each $A$ and $B$ is a free product of at least two cyclic groups of order $2$, and \eqref{eq_fo} implies that $x_1\ldots x_r$ and $x_{r+1}\ldots x_n$ do not have finite order in $A$ and in $B$, respectively. 
 
 By Lemma \ref{lem: rec presented}, it suffices to show that $L_n$ has solvable word problem. Note that $A,B$ have solvable word problem by an argument similar to the one in Lemma~\ref{lem Kn}. We  now   proceed via the  usual normal form of elements in an amalgam of two groups.  
 Let $U = A\cap B$ and note  \[U=\langle x_1\ldots x_r\rangle = \langle x_{r+1}\ldots x_n\rangle.\] Below we will pick suitable right coset representatives $\mathcal{S}$ of $U$ in $A$,  and $\mathcal{T}$ of $U$ in $B$, both containing~$1$. We first show that each element of $L_n$ can be uniquely written as $u s_1t_1 \ldots s_\ell t_\ell$ where $u\in U$ and each $s_i\in \mathcal{S}$ and $t_i\in \mathcal{T}$, and the $s_i, t_j$ are non-trivial with the possible exceptions $s_1$ and $t_\ell$:
  note that every element in $ A*_\ZZ B$ has the form $a_1b_1\ldots a_\ell b_\ell$ with each $a_i\in A$ and $b_i\in B$; starting from the right, replace $b_\ell = ub'_\ell$ with $u\in U$ and $b'_\ell\in \mathcal{T}$. Recall that $u\in A\cap B$, so next we replace $a_\ell u = va_\ell'$ with $v\in U$ and $a_\ell'\in \mathcal{S}$. An iteration of this process yields the required form. 

In order to   write an algorithm to create and multiply normal forms, we need to discuss the computability of $\mathcal{S}$; the discussion for $\mathcal{T}$ will be similar.  Recall that a  normal form for elements of $A$ is given by words $x_{k_1} \ldots x_{k_m}$ where each $k_i \in \{1, \ldots, r\}$ and $k_i \neq k_{i+1}$. Dropping the $x$'s, we can describe this as a sequence $\sss=[k_1, \ldots, k_m]$ of numbers in $\{1, \ldots, r\}$ as above. Multiplication of two sequences is induced by multiplication by elements in $A$, with the cancellation rules $x_i^2=1$ for each $i$. Thus, given two such sequences $\sigma$ and $\tau$, one can write $\sigma=\sigma'\varrho$ and $\tau=\varrho^{-1}\tau'$ where the subsequence $\varrho$ is as long as possible; in this case, $\sigma\tau=\sigma'\tau'$. (E.g.\ if $\sigma=[4,2,1]$ and $\tau=[1,2,1,3]$, then $\varrho=[2,1]$ and  $\sigma\tau=[4,1,3]$.) In the following we  use the notion of computable sets, implicitly assuming an encoding of group elements by natural numbers according to Remark~\ref{Malcev}. The subgroup $U = \la x_1 \ldots x_r \ra$ of $A$  is   computable, because given a normal form for $A$ one can check whether it  represents a power of $x_1 \ldots x_r$. As   a right transversal $\mathcal{S}$ of $U$ in $A$, we now pick those elements $\sss$ that are the length-lexicographically least in their right cosets of $\la x_1 \ldots x_r\ra  $. The latter is decidable because $\mathcal{S}$ consists of  all the sequences $\sss$ that do not start with $[1, 2,\ldots, s]$ or with $[r,r-1,\ldots,r-t]$ for some $s> r/2$ or $t< r/2$: to see this, observe that  \[\langle x_1\ldots x_r\rangle x_1\ldots x_s = \langle x_1\ldots x_r\rangle x_r x_{r-1}\ldots x_{s+1},\] and $[r,r-1,\ldots,s+1]$ is lexicographically smaller than $[1,2,\ldots,s]$; analogously for the representative $x_rx_{r-1}\ldots x_{r-t}$.  Since $r$ is fixed, $\mathcal{S}$ is computable. A similar argument for $\mathcal{T}$ (with variables $y_i = x_{n-i+1}$ for $1\le i \le n-r$) shows that $\mathcal{T}$ is computable. Thus $L_n$ has solvable word problem, and  so does  $H_n$.
\end{proof}

\end{document}
